% =====================================================================
%  整理者　macanine <macanine@163.com>
%  仓库　　github.com/macanine/zju-linear-algebra-papers
%  来源　　解析据 papers/2022-2023-秋冬-期中-甲.tex 撰写（原卷见 sources/2016-2025-期中-甲-历年真题汇编.pdf 第 11 页）
% =====================================================================

\begin{solution}
各行的四个数之和都是 $2+4+6+8=20$，故把第 $2,3,4$ 列统统加到第 $1$ 列，第 $1$ 列的元素全变成 $20$，提出公因子 $20$：
\[
D_1=20\begin{vmatrix}1&4&6&8\\1&6&8&2\\1&8&2&4\\1&2&4&6\end{vmatrix}.
\]
再用行倍加 $R_i\leftarrow R_i-R_1$（$i=2,3,4$）把第 $1$ 列的后三个元素消成零，并按第 $1$ 列展开：
\[
D_1=20\begin{vmatrix}1&4&6&8\\0&2&2&-6\\0&4&-4&-4\\0&-2&-2&-2\end{vmatrix}
=20\begin{vmatrix}2&2&-6\\4&-4&-4\\-2&-2&-2\end{vmatrix}.
\]
右边三阶行列式按第一行展开，两个二阶余子式分别是
$\begin{vmatrix}-4&-4\\-2&-2\end{vmatrix}=0$ 与 $\begin{vmatrix}4&-4\\-2&-2\end{vmatrix}=-16$，于是
\[
D_1=20\Bigl(2\cdot0-2\cdot(-16)+(-6)\cdot(-16)\Bigr)=20\cdot128=2560 .
\]
\ans{$D_1=2560$。}
\end{solution}

\begin{solution}
\textbf{构造新行列式。} 记 $D'$ 为把 $D$ 的第一行换成 $(0,2^{2},3^{2},\cdots,(n+1)^{2})$（即第 $i+1$ 列放 $(i+1)^{2}$）、其余各行不动所得到的行列式。$D$ 与 $D'$ 只在第一行不同，而代数余子式只依赖于划去该行该列后剩下的子式，所以 $D$ 中 $x_i$ 所在位置（第一行第 $i+1$ 列）的代数余子式在 $D'$ 中仍然是 $A_i$。把 $D'$ 按第一行展开，第一列的元素是 $0$，于是
\[
D'=0\cdot A_0+2^{2}A_1+3^{2}A_2+\cdots+(n+1)^{2}A_n=\sum_{i=1}^{n}(i+1)^{2}A_i,
\]
（$A_0$ 是第一行第一列元素的代数余子式，它乘 $0$，不影响结果。）这说明所求的组合恰好就是 $D'$ 的值。

\textbf{计算 $D'$。} $D'$ 的第一列是 $(0,1,1,\cdots,1)^{T}$，按第一列展开得
\[
D'=\sum_{j=2}^{n+1}(-1)^{j+1}M_{j1},
\]
$M_{j1}$ 是划去 $D'$ 的第 $j$ 行与第 $1$ 列后余下的 $n$ 阶子式。划去后，原第一行剩下 $(2^{2},3^{2},\cdots,(n+1)^{2})$，而原来的第 $\ell$ 行（$\ell=2,\cdots,n+1$，$\ell\ne j$）只剩下一个非零元：数值为 $\ell$、位于第 $\ell$ 列。把 $(2^{2},3^{2},\cdots,(n+1)^{2})$ 这一行逐次与相邻行交换、下移到第 $j-1$ 行（共 $j-2$ 次交换，代价符号 $(-1)^{j-2}$）。此时前 $j-2$ 行各只有一个非零元 $2,3,\cdots,j-1$，沿它们依次展开后余下的子块为上三角阵，其对角元依次为 $j^{2},j+1,\cdots,n+1$，于是
\[
M_{j1}=(-1)^{j-2}\cdot j^{2}\cdot\prod_{\substack{\ell=2\\ \ell\ne j}}^{n+1}\ell
=(-1)^{j-2}j^{2}\cdot\frac{(n+1)!}{j}=(-1)^{j-2}j\,(n+1)!,
\]
最后一步用了 $\displaystyle\prod_{\ell=2}^{n+1}\ell=(n+1)!$。代入 $(-1)^{j+1}(-1)^{j-2}=(-1)^{2j-1}=-1$，得
\[
D'=-(n+1)!\sum_{j=2}^{n+1}j=-(n+1)!\cdot\frac{(n+1)(n+2)-2}{2}=-\frac{n(n+3)}{2}(n+1)! .
\]
\ans{$4A_1+9A_2+\cdots+(n+1)^{2}A_n=-\dfrac{n(n+3)}{2}(n+1)!$。}

验算：取 $n=2$ 时 $D$ 为三阶行列式，直接算得 $A_1=-3$，$A_2=-2$，所求为 $4\times(-3)+9\times(-2)=-30$，与公式 $-\dfrac{2\times5}{2}\cdot3!=-30$ 一致。
\end{solution}

\begin{solution}
对增广矩阵作初等行变换：先交换第 $1,3$ 行，再作行倍加 $R_2-R_1$ 与 $R_3-\lambda R_1$，得
\[
\left(\begin{array}{cccc|c}
1&1&\lambda&1&\lambda^{2}\\
0&\lambda-1&1-\lambda&0&\lambda-\lambda^{2}\\
0&1-\lambda&1-\lambda^{2}&1-\lambda&1-\lambda^{3}
\end{array}\right).
\]

\textbf{先设 $\lambda\ne1$}，此时可把后两行分别除以 $\lambda-1$ 与 $1-\lambda$（用 $1-\lambda^{3}=(1-\lambda)(1+\lambda+\lambda^{2})$），再作 $R_3-R_2$：
\[
\left(\begin{array}{cccc|c}
1&1&\lambda&1&\lambda^{2}\\
0&1&-1&0&-\lambda\\
0&0&\lambda+2&1&(\lambda+1)^{2}
\end{array}\right).
\]
当 $\lambda\ne-2$ 时系数矩阵的秩为 $3$，取 $x_4=t$ 为自由未知量，回代得
\[
x_3=\frac{(\lambda+1)^{2}-t}{\lambda+2},\qquad
x_2=x_3-\lambda=\frac{1-t}{\lambda+2},\qquad
x_1=\lambda^{2}-x_2-\lambda x_3-x_4=-\frac{\lambda+1+t}{\lambda+2},
\]
即通解
\[
(x_1,x_2,x_3,x_4)^{T}
=\frac{1}{\lambda+2}\bigl(-(\lambda+1),\,1,\,(\lambda+1)^{2},\,0\bigr)^{T}
+t\cdot\frac{1}{\lambda+2}\bigl(-1,\,-1,\,-1,\,\lambda+2\bigr)^{T}\qquad(t\in\mathbb{R}).
\]
当 $\lambda=-2$ 时第三行变成 $x_4=(\lambda+1)^{2}=1$，系数矩阵的秩仍为 $3$，取 $x_3=t$，由 $x_2-x_3=-\lambda=2$ 得
\[
(x_1,x_2,x_3,x_4)^{T}=(1,2,0,1)^{T}+t\,(1,1,1,0)^{T}\qquad(t\in\mathbb{R}).
\]

\textbf{再设 $\lambda=1$}，三个方程都成为 $x_1+x_2+x_3+x_4=1$，系数矩阵的秩为 $1$，解集是一个三维平面：
\[
x_1=1-t_1-t_2-t_3,\quad x_2=t_1,\quad x_3=t_2,\quad x_4=t_3\qquad(t_1,t_2,t_3\in\mathbb{R}).
\]
\ans{方程组总有无穷多解：$\lambda\ne1,-2$ 时通解为
$\bigl(-(\lambda+1),1,(\lambda+1)^{2},0\bigr)^{T}/(\lambda+2)
+t\bigl(-1,-1,-1,\lambda+2\bigr)^{T}/(\lambda+2)$；
$\lambda=-2$ 时通解为 $(1,2,0,1)^{T}+t(1,1,1,0)^{T}$；
$\lambda=1$ 时解集为 $x_1+x_2+x_3+x_4=1$。}
\end{solution}

\begin{solution}
把第 $2,3,4$ 行分别减去第 $1$ 行的 $1,3,5$ 倍：
\[
A\to\begin{pmatrix}1&2&5&-1&-1\\0&0&-2&2&4\\0&0&-2&2&4\\0&0&-4&4&8\end{pmatrix}.
\]
此时第 $3$ 行与第 $2$ 行相同，第 $4$ 行是第 $2$ 行的 $2$ 倍，故再作 $R_3-R_2$、$R_4-2R_2$ 得
\[
\to\begin{pmatrix}1&2&5&-1&-1\\0&0&-2&2&4\\0&0&0&0&0\\0&0&0&0&0\end{pmatrix}
\to\begin{pmatrix}1&2&5&-1&-1\\0&0&1&-1&-2\\0&0&0&0&0\\0&0&0&0&0\end{pmatrix}=B,
\]
最后一步是把第 $2$ 行乘以 $-\dfrac12$。$B$ 的阶梯头（每行第一个非零元）分别为第 $1,3$ 列的 $1$，已是所要求的行阶梯形，且 $r(A)=r(B)=2$。若继续把第 $1$ 行减去 $5$ 倍的第 $2$ 行，可得行最简形
$\begin{pmatrix}1&2&0&4&9\\0&0&1&-1&-2\\0&0&0&0&0\\0&0&0&0&0\end{pmatrix}$。
\ans{$B=\begin{pmatrix}1&2&5&-1&-1\\0&0&1&-1&-2\\0&0&0&0&0\\0&0&0&0&0\end{pmatrix}$（阶梯头在第 $1,3$ 列，$r(A)=r(B)=2$）。}
\end{solution}

\begin{solution}
\textbf{(1) 秩的定义。} 设 $A$ 是 $m\times n$ 矩阵。如果在 $A$ 中有一个 $r$ 阶子式不为零，而所有 $r+1$ 阶子式（如果存在的话）全为零，就称 $A$ 的秩为 $r$，记作 $r(A)$；即 $r(A)$ 就是 $A$ 中非零子式的最高阶数。零矩阵没有非零子式，规定 $r(O)=0$。（由此 $0\le r(A)\le\min\{m,n\}$。）

\textbf{(2)} \emph{必要性}（$r(A)\le r$ $\Rightarrow$ 所有 $r+1$ 阶子式为零）：若 $A$ 有一个非零的 $r+1$ 阶子式，则 $A$ 中非零子式的最高阶数至少为 $r+1$，即 $r(A)\ge r+1$，矛盾。

\emph{充分性}（所有 $r+1$ 阶子式为零 $\Rightarrow$ $r(A)\le r$）：先用一个事实——若非零的 $s$ 阶子式 $D$ 存在，则 $D$ 中必有一行不全为零，把 $D$ 按这一行展开，$D$ 等于该行各元素与其代数余子式乘积之和；既然 $D\ne0$，至少有一个 $s-1$ 阶余子式不为零，即 $A$ 中存在非零的 $s-1$ 阶子式。假若 $r(A)=s>r$，取一个非零的 $s$ 阶子式，反复使用上述事实，便从非零的 $s$ 阶子式依次得到非零的 $s-1$ 阶、$s-2$ 阶、$\cdots$、$r+1$ 阶子式，与“所有 $r+1$ 阶子式为零”矛盾，故 $r(A)\le r$。（当 $r+1>\min\{m,n\}$ 时 $A$ 没有 $r+1$ 阶子式，条件空真，而此时 $r(A)\le\min\{m,n\}\le r$，结论同样成立。）

\textbf{(3)} 设 $B$ 由 $A$ 经一次行倍加得到，即 $B$ 的第 $j$ 行等于 $A$ 的第 $j$ 行加上第 $i$ 行的 $k$ 倍（$i\ne j$），其余各行与 $A$ 的对应行相同。

考察 $B$ 的任一 $s$ 阶子式，设它的行取自 $B$ 的第 $r_1<\cdots<r_s$ 行。若行标中不含 $j$，该子式就是 $A$ 的同一个 $s$ 阶子式；若含 $j$，则由行列式对行的线性性，该子式的行列式等于 $D_1+kD_2$，其中 $D_1$ 是把这一行换回 $A$ 的第 $j$ 行所得的行列式，$D_2$ 是把这一行换成 $A$ 的第 $i$ 行所得的行列式：$D_1$ 是 $A$ 的 $s$ 阶子式；$D_2$ 当 $i$ 也在行标中时因两行相同而为 $0$，否则也是 $A$ 的 $s$ 阶子式。总之，$B$ 的每个 $s$ 阶子式都是 $A$ 的 $s$ 阶子式的线性组合。

取 $s=r(A)+1$：由 (2)，$A$ 的 $r(A)+1$ 阶子式全为零，故 $B$ 的 $r(A)+1$ 阶子式也全为零，再由 (2) 得 $r(B)\le r(A)$。反过来，$A$ 也可由 $B$ 经一次行倍加得到（第 $j$ 行加上第 $i$ 行的 $-k$ 倍），同理 $r(A)\le r(B)$。故 $r(B)=r(A)$。
\qed
\end{solution}

\begin{solution}
\textbf{(1)} 分三种情形。
若 $r(A)=n$：$A$ 可逆，由 $A^{*}=|A|A^{-1}$ 知 $A^{*}$ 也可逆，故 $r(A^{*})=n$。
若 $r(A)=n-1$：$A$ 有非零的 $n-1$ 阶子式，而 $A^{*}$ 的每个元素都是 $A$ 的 $n-1$ 阶代数余子式，故 $A^{*}\ne O$，$r(A^{*})\ge1$。另一方面 $|A|=0$，由基本等式 $AA^{*}=|A|E$ 得 $AA^{*}=O$，即 $A^{*}$ 的每一列 $X$ 都满足 $AX=\theta$，于是
\[
\operatorname{Col}(A^{*})\subseteq\ker A .
\]
由秩－零化度定理（并把秩看作列空间的维数）得 $\dim\ker A=n-r(A)=1$；又 $r(A^{*})=\dim\operatorname{Col}(A^{*})$，故 $r(A^{*})\le1$。两端合起来得 $r(A^{*})=1$。
若 $r(A)<n-1$：$A$ 的所有 $n-1$ 阶子式全为零，故 $A^{*}=O$，$r(A^{*})=0$。

\textbf{(2)} 当 $|A|\ne0$ 时 $A^{*}=|A|A^{-1}$。对 $A^{*}A=|A|E$ 两边取行列式：$|A^{*}|\cdot|A|=|A|^{n}$，故 $|A^{*}|=|A|^{n-1}$，于是
\[
(A^{*})^{*}=|A^{*}|(A^{*})^{-1}=|A|^{n-1}\cdot\frac{A}{|A|}=|A|^{n-2}A .
\]
当 $|A|=0$ 时 $r(A)\le n-1$，由 (1) 知 $r(A^{*})\le1$。若 $n\ge3$，则 $r(A^{*})\le1\le n-2<n-1$，把 (1) 用于 $n$ 阶矩阵 $A^{*}$ 得 $(A^{*})^{*}=O$；又此时 $n-2\ge1$，$|A|^{n-2}=0$，故仍成立 $(A^{*})^{*}=|A|^{n-2}A$。若 $n=2$，规定 $|A|^{0}=1$，直接验证（$A=\begin{pmatrix}a&b\\c&d\end{pmatrix}$，$A^{*}=\begin{pmatrix}d&-b\\-c&a\end{pmatrix}$）得 $(A^{*})^{*}=\begin{pmatrix}a&b\\c&d\end{pmatrix}=A=|A|^{0}A$。故对一切 $n\ge2$ 有 $(A^{*})^{*}=|A|^{n-2}A$。

\textbf{(3)} 由 (2)，$(A^{*})^{*}=|A|^{n-2}A$，于是
\[
(A^{*})^{*}=A\iff|A|^{n-2}A=A\iff\bigl(|A|^{n-2}-1\bigr)A=O .
\]
若 $|A|^{n-2}=1$，显然右边的等式成立，即 $(A^{*})^{*}=A$。反之设 $(|A|^{n-2}-1)A=O$：条件 $r(A)>0$ 等价于 $A\ne O$，取 $A$ 的一个非零元素 $a_{pq}$，比较矩阵等式两边的 $(p,q)$ 元得 $(|A|^{n-2}-1)a_{pq}=0$，故 $|A|^{n-2}=1$。因此当 $r(A)>0$ 时 $(A^{*})^{*}=A$ 当且仅当 $|A|^{n-2}=1$。
\qed
\end{solution}

\begin{solution}
先明确各矩阵的阶：要使 $AB$ 与 $CD$ 同时有意义，只能是 $A,C$ 为 $m\times n$ 阶矩阵、$B,D$ 为 $n\times s$ 阶矩阵，以下按此理解。

在乘积中间插入 $CB$，把差拆成两项：
\[
AB-CD=AB-CB+CB-CD=(A-C)B+C(B-D).
\]
用到两条基本性质（把秩看作列空间的维数即可）：
\[
r(X+Y)\le r(X)+r(Y),\qquad r(XY)\le\min\{r(X),r(Y)\}.
\]
第一条因为 $X+Y$ 的列空间含于 $\operatorname{Col}(X)+\operatorname{Col}(Y)$，而两个子空间之和的维数不超过维数之和；第二条因为 $XY$ 的每一列都是 $Y$ 的列向量的线性组合，故 $\operatorname{Col}(XY)\subseteq\operatorname{Col}(Y)$，对 $r(XY)\le r(X)$ 取转置即得。于是
\[
r(AB-CD)\le r\bigl((A-C)B\bigr)+r\bigl(C(B-D)\bigr)\le r(A-C)+r(B-D).
\]
\qed
\rem{题面“$A,C,D$ 均为 $m\times n$ 阶矩阵”中的 $D$ 是笔误：$D$ 不可能同时是 $m\times n$ 阶与 $n\times s$ 阶，应读作“$A,C$ 均为 $m\times n$ 阶矩阵”。}
\end{solution}

\begin{solution}
记 $d_k=r(A^{k})-r(A^{k+1})$（$k\ge1$）。

\textbf{第一步：相邻两次掉秩是递减的。} 用 Frobenius 秩不等式
\[
r(XY)+r(YZ)\le r(Y)+r(XYZ),
\]
取 $X=A$，$Y=A^{k}$，$Z=A$（$k\ge1$），得
\[
2r(A^{k+1})\le r(A^{k})+r(A^{k+2}),
\]
即 $d_k\ge d_{k+1}$。故 $d_1\ge d_2\ge\cdots$，特别地 $d_1\ge d_{n+2}$，也就是
\[
r(A^{n+2})+r(A^{2})\le r(A^{n+3})+r(A). \qquad(\star)
\]

\textbf{第二步：对 $n$ 归纳，证明 $(n+1)r(A^{2})\le r(A^{n+2})+n\,r(A)$。}
$n=1$ 时即 $2r(A^{2})\le r(A^{3})+r(A)$，这正是 $d_1\ge d_2$。
设 $n$ 时结论成立，两边同时加上 $r(A^{2})$：
\[
(n+2)r(A^{2})\le r(A^{n+2})+r(A^{2})+n\,r(A)\le r(A^{n+3})+r(A)+n\,r(A)=r(A^{n+3})+(n+1)r(A),
\]
中间一步用的正是 $(\star)$。故结论对 $n+1$ 也成立。

\textbf{第三步：把 $m$ 放进去。} 由 $r(A)\ge0$，对任何 $m\ge n$ 都有
\[
(n+1)r(A^{2})\le r(A^{n+2})+n\,r(A)\le r(A^{n+2})+m\,r(A),
\]
这就是所要证的不等式；$m$ 取得更大时右端更大，不等式自然仍成立。（证明与矩阵的阶无关，$2022^{2022}$ 只是题面给定的方阵阶数。）
\qed
\rem{题面写“$\forall n\in Z^{+},\ \forall m\in Z^{+}$”，但按字面 $m<n$ 时结论不成立：取 $A$ 为 $N=2022^{2022}$ 阶单位阵 $E$，则 $r(A^{k})=N$，不等式成为 $(n+1)N\le(1+m)N$，只在 $m\ge n$ 时成立。上面的归纳证得的正是系数取 $n$ 的最强形式 $(n+1)r(A^{2})\le r(A^{n+2})+n\,r(A)$，故题面中的 $m$ 应理解为 $m\ge n$。}
\end{solution}
